\documentclass[11pt]{article}
\usepackage[margin=1in]{geometry}
\usepackage{hyperref}
\usepackage{enumitem}
\usepackage{booktabs}
\usepackage{xcolor}
\title{ProbOS --- Month 1, Week 2\\Model ABC + BatteryModel2Cell}
\author{Nisong Monyimba}
\date{\today}

\begin{document}
\maketitle

\section*{Week 2 goal}
Build the \texttt{Model} abstract base class defining the
\texttt{forward\_batch()} contract every domain model must satisfy, and
the first concrete model: \texttt{BatteryModel2Cell}, an 8-state
Arrhenius thermal-runaway ODE validated against Kim et al. (2007).

\section*{Day-by-day plan}

\subsection*{Day 1 -- Model ABC design}
\begin{itemize}[leftmargin=*]
  \item Define the interface: \texttt{state\_dim}, \texttt{param\_dim},
    \texttt{param\_names()}, \texttt{initial\_state()},
    \texttt{forward\_batch(state, params, dt)}
  \item Key constraint: \texttt{forward\_batch} must be vectorised over
    $N$ particles via NumPy broadcasting, no per-particle Python loops
  \item \texttt{python/src/state.py} skeleton
\end{itemize}

\subsection*{Day 2 -- BatteryModel2Cell state and parameters}
\begin{itemize}[leftmargin=*]
  \item Define the 8-state vector: two cells $\times$ (temperature,
    SEI concentration, anode concentration, cathode concentration)
  \item Define the 15-parameter vector: activation energies,
    pre-exponential factors, heat capacities per reaction
  \item \texttt{python/src/battery\_model.py} skeleton
\end{itemize}

\subsection*{Day 3 -- Arrhenius kinetics implementation}
\begin{itemize}[leftmargin=*]
  \item Implement the three parallel exothermic reactions (SEI
    decomposition, anode reaction, cathode reaction) via the Arrhenius
    rate law
  \item Explicit Euler integration of the coupled ODE system
  \item Vectorise fully across particles
\end{itemize}

\subsection*{Day 4 -- Parameter priors}
\begin{itemize}[leftmargin=*]
  \item \texttt{python/src/parameter\_priors.py}: 15 prior distributions
    built from the Week 1 \texttt{Distribution} types, centred on Kim
    et al. (2007) nominal values
  \item \texttt{build\_battery\_priors()} factory function
\end{itemize}

\subsection*{Day 5 -- Kim 2007 validation}
\begin{itemize}[leftmargin=*]
  \item Deterministic run at nominal parameters, adiabatic conditions
  \item Compare onset temperature and thermal runaway timing against
    Kim et al. (2007) accelerating rate calorimetry (ARC) data
  \item \texttt{python/examples/week2\_battery\_ode.py} deterministic
    demo script
\end{itemize}

\subsection*{Day 6 -- CLT convergence demo}
\begin{itemize}[leftmargin=*]
  \item \texttt{python/examples/week2\_clt\_demo.py}: verify Monte
    Carlo error shrinks at the theoretical $1/\sqrt{N}$ rate on a
    simple test case, ahead of building the full engine in Week 3
\end{itemize}

\subsection*{Day 7 -- Test coverage + hardening}
\begin{itemize}[leftmargin=*]
  \item \texttt{python/tests/test\_state.py} -- Model ABC contract
    tests
  \item \texttt{python/tests/test\_battery\_model.py} -- 8-state ODE
    correctness, Kim 2007 validation as an automated test, physical
    invariants (temperatures stay positive, concentrations stay in
    $[0,1]$)
  \item mypy strict, ruff clean, CI green
\end{itemize}

\section*{What was actually accomplished (retrospective)}

\begin{itemize}[leftmargin=*]
  \item \texttt{Model} ABC implemented exactly as planned
  \item \texttt{BatteryModel2Cell}: 8-state Arrhenius thermal-runaway
    ODE, fully vectorised \texttt{forward\_batch()}
  \item 15 prior distributions in \texttt{parameter\_priors.py}, built
    entirely from Week 1's \texttt{Distribution} types -- no new
    distribution code needed
  \item Kim et al. (2007) validation: onset temperature and runaway
    behaviour match the reference ARC data
  \item Deterministic ODE demo and CLT convergence demo scripts
  \item 67 new Python tests (test\_state.py + test\_battery\_model.py)
\end{itemize}

\section*{Numbers at Week 2 close}
\begin{itemize}[leftmargin=*]
  \item Python tests: 111/111 passing (44 Week 1 + 67 new)
  \item C++ tests: 13/13 passing (unchanged from Week 1)
  \item mypy strict: 0 errors
  \item ruff: 0 warnings
\end{itemize}

\section*{Carried into Week 3}
\texttt{BatteryModel2Cell} + \texttt{build\_battery\_priors()} become
the standard test case for every subsequent kernel layer: the Monte
Carlo engine, Sobol sensitivity analysis, and provenance tracking in
Week 3 are all built and validated against this exact model.

\end{document}
